\section{Evidence and Provenance Boundary}
\label{app:evidence-provenance}

The maintained evidence is retained under:
\begin{itemize}
 \item the direct unsealed Mathieu1D execution bound by
       \path{experiment-references.json};
 \item \path{calculations/ICMSEP2026/}\\
       \path{conference/paper_1/multiband-alignment/};
 \item \path{calculations/ICMSEP2026/}\\
       \path{conference/paper_1/constrained-admissible-sets/};
 \item \path{calculations/ICMSEP2026/}\\
       \path{conference/paper_1/}\\
       \path{constrained-admissible-sets-threshold-sensitivity/};
 \item \path{calculations/ICMSEP2026/}\\
       \path{conference/paper_1/}\\
       \path{silicon-wannier-kinetic-preview/};
 \item \path{calculations/research-monograph/periodic-2d/}; and
 \item \path{calculations/ICMSEP2026/conference/paper_1/}.
\end{itemize}
The retained directories contain, as applicable, protocols, result records,
independent verification code, figures, reports, source-correlation records, and
checksum manifests. The direct Mathieu1D execution, M2--M3, and two-dimensional evidence is controlled
illustrative software and numerical verification. The silicon directory is a
bounded extraction from a read-only analysis of official upstream tutorial
inputs. None is a public data deposit, DOI, material validation, or uncertainty
analysis.

The publication-independent direct Mathieu1D execution contains the retained lowest-band
calculation, sampled external-gap diagnostics, rank/conditioning channels, typed
constructor contracts, and byte-owning independent verification. It is explicitly
unsealed. Its sampled gaps do not establish continuum isolation, and no predecessor
dataset remains a live input.
The multiband package contains the prospectively frozen
M2 rank-two alignment and locality calculation, including its independent
reconstruction and disjoint withheld diagnostics. The constrained-admissible-
set package contains M3's prospectively frozen two-parameter, nine-component
finite alignment family, common witness, and bounded separation certificate.
The threshold-sensitivity package is a separate post-hoc analytic reanalysis of
the sealed M3 quadratics and is not part of the confirmatory package. The
silicon preview copies a small content-bound subset of an external analysis and
adds a post-hoc correlation confirming that all four original source-run files
match the official Wannier90~3.1.0 \path{examples/example11} tutorial by
SHA-256. The executed \path{silicon.win} adds only two bound output-artifact
requests and changes no tutorial numerical setting. The preview supports
tutorial-scale same-frame construction and finite-mesh truncation
diagnostics only. It does not authenticate the causal output lineage,
establish convergence, or report effective masses. The two-dimensional reports
contain additional composite and topology controls.
The conference-paper parent calculation contains the direct quadratic
admissible-set map, exact common witness, rational separation certificate,
independent verifier, and figure inputs. Only the bounded results
directly used in the paper are summarized in its main text. Historical external
executions are not rerun by the manuscript.

\subsection{Verifier and digest semantics}

The direct Mathieu1D byte-owning verifier and the standalone M2 and M3 retained
verifiers reconstruct their finite protocols without calling the corresponding
producer Actions. The direct verifier checks canonical typed bytes,
constructor-restored correlations and exact subject and projection digests. The separate
threshold-sensitivity verifier
reconstructs its analytic grid after checking both the source-result digest and
that result's exact entry in the source package manifest. These verifiers
remain dependent on shared numerical libraries, runtime behavior, and declared
mathematical conventions; passing establishes bounded reconstruction
consistency rather than an independent physical oracle.
The checksum manifests bind retained bytes to recorded SHA-256 digests and
therefore detect inconsistency relative to those manifests. A digest manifest
alone is not an external trusted timestamp and does not, without an independent
commit or archival record, prove when the bytes were created.

\subsection{Curation work}

The remaining evidence-curation tasks are:
\begin{itemize}
 \item[$\square$] trace every numerical statement and plotted value in the
 paper to a retained result field or report table;
 \item[$\square$] verify the figure sources against the retained checksum
 manifests;
 \item[$\square$] record the exact software revision and runtime environment
 associated with each retained calculation;
 \item[$\square$] decide which machine-readable records and figure-generation
 inputs belong in a public archival package; and
 \item[$\square$] obtain authorization before publication or deposition.
\end{itemize}
These practices follow general reproducible-computation guidance~\cite{sandve2013}.
Completing the tasks strengthens provenance and publication readiness; it does
not change the evidence class from controlled numerical verification to
material validation.
